% ============================================================
%  EXAMEN TEMA 3: FUNCIONES LINEALES
%  Curso de Introducción — 2do Año
%  Duración: 90 minutos
%  Temas: Conceptos básicos, plano cartesiano, tipo de función,
%         gráficas, pendiente y modelado
%  Autor: Ing. Gabriel Astudillo
%  Nota: enunciado limpio (sin pistas ni desarrollo). Las
%  soluciones están en la "Clave de Respuestas" al final.
% ============================================================

\documentclass[12pt, letterpaper]{article}

% ── Paquetes ─────────────────────────────────────────────────

\usepackage{mathwizards-guia}
\mwguia{Examen — Tema 3: Funciones Lineales}{Ing. Gabriel Astudillo $\cdot$ Curso 2do Año}

\begin{document}
% ============================================================

% ── Portada / Título ─────────────────────────────────────────
\mwportada{Examen del Tema 3}{Funciones Lineales}{Función Lineal y Afín, Plano Cartesiano, Gráficas, Pendiente y Modelado}

\vspace{4pt}
\begin{cuadroinfo}
  \small
  \textbf{Nombre:} \rule{0.55\linewidth}{0.4pt} \quad \textbf{Fecha:} \rule{0.15\linewidth}{0.4pt} \\[6pt]
  \textbf{Tiempo disponible:} 90 minutos \quad$\bullet$\quad \textbf{Puntaje total:} 100 puntos \\[4pt]
  \textbf{Instrucciones:}
  \begin{enumerate}[leftmargin=*]
    \item Lee cada ejercicio con calma antes de resolverlo.
    \item Muestra \emph{todos} tus pasos y traza las gráficas con regla.
    \item Recuerda: \emph{lineal} es $F(x) = mx$ (pasa por el origen); \emph{afín} es $F(x) = mx + b$ con $b \neq 0$.
    \item En los problemas de aplicación, indica claramente qué representa cada variable.
    \item No se permite el uso de calculadora.
    \item Escribe con letra clara y revisa tu examen antes de entregarlo.
  \end{enumerate}
\end{cuadroinfo}

\vspace{8pt}

% ============================================================
\section{Conceptos Básicos: Función Lineal y Afín (15 puntos)}
% ============================================================

\begin{enumerate}[leftmargin=*]
  \item \textbf{(5 pts)} Clasifica cada función en lineal (proporcionalidad directa), afín o constante, e indica su pendiente $m$ y su coeficiente de posición $b$:
        \[
          a)\ F(x) = -7x + 2 \qquad b)\ F(x) = \frac{3}{4}x \qquad c)\ F(x) = -9 \qquad d)\ F(x) = 5 - x \qquad e)\ F(x) = 0x + 4
        \]

  \item \textbf{(5 pts)} Si $F(x) = -4x + 3$, calcula:
        \[
          F(0), \qquad F(2), \qquad F(-3), \qquad F\!\left(\frac{1}{2}\right)
        \]

  \item \textbf{(5 pts)} Si $F(x) = mx + b$, con $F(2) = 7$ y $F(5) = 16$, determina $m$ y $b$.
\end{enumerate}

\newpage

% ============================================================
\section{Plano Cartesiano y Coordenadas (12 puntos)}
% ============================================================

\begin{enumerate}[leftmargin=*]
  \item \textbf{(4 pts)} Indica el cuadrante o el eje donde se ubica cada punto:
        \[
          a)\ (-6, 2) \qquad b)\ (5, -3) \qquad c)\ (0, -7) \qquad d)\ (-4, 0) \qquad e)\ (-3, -8)
        \]

  \item \textbf{(4 pts)} Escribe las coordenadas del punto descrito:
        \begin{enumerate}[label=\alph*)]
          \item $5$ unidades a la izquierda y $4$ unidades arriba del origen.
          \item $2$ unidades a la derecha y $6$ unidades abajo del origen.
        \end{enumerate}

  \item \textbf{(4 pts)} Un punto está en el cuadrante II y sus coordenadas son números enteros cuya suma es $3$. Escribe dos posibles coordenadas para ese punto.
\end{enumerate}

\newpage

% ============================================================
\section{Identificación del Tipo de Función (15 puntos)}
% ============================================================

\begin{enumerate}[leftmargin=*]
  \item \textbf{(5 pts)} Determina si la función que pasa por los puntos dados es creciente, decreciente o constante:
        \[
          a)\ (0,5) \text{ y } (3,5) \qquad b)\ (1,2) \text{ y } (4,14) \qquad c)\ (2,9) \text{ y } (5,3)
        \]

  \item \textbf{(5 pts)} Halla la pendiente y el coeficiente de posición de la recta que pasa por los puntos $(2,3)$ y $(6,11)$. Escribe su ecuación.

  \item \textbf{(5 pts)} La recta pasa por los puntos $(4,-2)$ y $(7,-11)$. Determina si es creciente o decreciente y escribe su ecuación.
\end{enumerate}

\newpage

% ============================================================
\section{Gráfica de Funciones Lineales (18 puntos)}
% ============================================================

\begin{enumerate}[leftmargin=*]
  \item \textbf{(6 pts)} Grafica por \textbf{tabulación} la función $F(x) = -2x + 4$, usando $x = -1, 0, 1, 2, 3$. Indica el tipo de función.

  \item \textbf{(6 pts)} Grafica por el \textbf{método de los puntos de corte}:
        \[
          a)\ F(x) = 2x - 6 \qquad b)\ F(x) = -\frac{3}{2}x + 6
        \]

  \item \textbf{(6 pts)} Sin graficar, determina los cortes con los ejes $X$ e $Y$ de $F(x) = \dfrac{4}{5}x + 8$ y clasifica la función (lineal, afín o constante).
\end{enumerate}

\newpage

% ============================================================
\section{Análisis de la Pendiente (20 puntos)}
% ============================================================

\begin{enumerate}[leftmargin=*]
  \item \textbf{(5 pts)} Calcula la pendiente de la recta que pasa por cada par de puntos:
        \[
          a)\ (-2,1) \text{ y } (3,11) \qquad b)\ (1,7) \text{ y } (4,-2) \qquad c)\ (3,5) \text{ y } (3,-4)
        \]

  \item \textbf{(5 pts)} Indica si cada función es creciente, decreciente o constante:
        \[
          a)\ F(x) = 7 - 3x \qquad b)\ F(x) = 2 \qquad c)\ F(x) = -x \qquad d)\ F(x) = \frac{5}{2}x
        \]

  \item \textbf{(5 pts)} Escribe la función afín cuya pendiente es $m = -4$ y que pasa por el punto $(2,-1)$.

  \item \textbf{(5 pts)} Una recta pasa por los puntos $(k, 6)$ y $(k+3, -3)$. Calcula su pendiente sin conocer el valor de $k$.
\end{enumerate}

\newpage

% ============================================================
\section{Problemas de Aplicación (Modelado) (20 puntos)}
% ============================================================

\begin{enumerate}[leftmargin=*]
  \item \textbf{(5 pts)} $4$ kg de naranjas cuestan $\$5\,600$. Escribe la función del costo en función de los kilos y calcula cuánto cuestan $7$ kg.

  \item \textbf{(5 pts)} Un taxi cobra $\$2\,000$ de bajada de bandera y $\$700$ por kilómetro. Escribe la función del costo y calcula el costo de un viaje de $9$ km.

  \item \textbf{(5 pts)} Un gimnasio cobra $\$12\,000$ de inscripción y $\$4\,000$ mensuales.
        \begin{enumerate}[label=\alph*)]
          \item Escribe la función del costo total en función de los meses.
          \item ¿Cuánto cuesta entrenar $8$ meses?
          \item ¿Cuántos meses se pueden pagar con $\$60\,000$?
        \end{enumerate}

  \item \textbf{(5 pts)} Un vendedor gana $\$300\,000$ fijos más el $8\%$ de sus ventas. Escribe la función de su sueldo en función de las ventas y calcula cuánto debe vender para ganar $\$460\,000$.
\end{enumerate}

\newpage

% ============================================================
\ifmwclaves
  \section{Clave de Respuestas}
  % ============================================================

  \subsection*{Sección 1: Conceptos Básicos}

  \begin{enumerate}[leftmargin=*, label=\textbf{\arabic*.}]
    \item a) afín, $m = -7$, $b = 2$ \quad b) lineal, $m = \dfrac{3}{4}$, $b = 0$ \quad c) constante, $m = 0$, $b = -9$ \quad d) afín, $m = -1$, $b = 5$ \quad e) constante, $m = 0$, $b = 4$

    \item $F(0) = 3$; \quad $F(2) = -4(2) + 3 = -5$; \quad $F(-3) = -4(-3) + 3 = 15$; \quad $F\!\left(\dfrac{1}{2}\right) = -4 \cdot \dfrac{1}{2} + 3 = 1$.

    \item Con $F(2) = 2m + b = 7$ y $F(5) = 5m + b = 16$. Restando: $3m = 9 \Rightarrow m = 3$.
          Luego $2(3) + b = 7 \Rightarrow b = 1$. Entonces $F(x) = 3x + 1$.
  \end{enumerate}

  \newpage

  \subsection*{Sección 2: Plano Cartesiano}

  \begin{enumerate}[leftmargin=*, label=\textbf{\arabic*.}]
    \item a) II \quad b) IV \quad c) eje $Y$ \quad d) eje $X$ \quad e) III

    \item a) $(-5, 4)$ \quad b) $(2, -6)$

    \item Deben cumplir $x < 0$, $y > 0$ y $x + y = 3$. Ejemplos: $(-1, 4)$ y $(-2, 5)$.
  \end{enumerate}

  \newpage

  \subsection*{Sección 3: Identificación del Tipo de Función}

  \begin{enumerate}[leftmargin=*, label=\textbf{\arabic*.}]
    \item a) constante ($y = 5$ en ambos) \quad b) creciente ($2 \to 14$) \quad c) decreciente ($9 \to 3$)

    \item $m = \dfrac{11 - 3}{6 - 2} = \dfrac{8}{4} = 2$; \quad $F(2) = 2(2) + b = 3 \Rightarrow b = -1$. \quad Ecuación: $F(x) = 2x - 1$.

    \item $m = \dfrac{-11 - (-2)}{7 - 4} = \dfrac{-9}{3} = -3 < 0 \Rightarrow$ \textbf{decreciente}.
          $F(4) = -3(4) + b = -2 \Rightarrow b = 10$. \quad Ecuación: $F(x) = -3x + 10$.
  \end{enumerate}

  \newpage

  \subsection*{Sección 4: Gráfica de Funciones Lineales}

  \begin{enumerate}[leftmargin=*, label=\textbf{\arabic*.}]
    \item Tabla de valores:
          \begin{center}
            \begin{tabular}{c|ccccc}
              \toprule
              $x$              & $-1$ & $0$ & $1$ & $2$ & $3$ \\
              \midrule
              $F(x) = -2x + 4$ & $6$  & $4$ & $2$ & $0$ & $-2$ \\
              \bottomrule
            \end{tabular}
          \end{center}
          Recta decreciente que pasa por $(0,4)$. Es una función \textbf{afín} (decreciente).

    \item a) $F(x) = 2x - 6$: corte con $Y$ en $F(0) = -6 \Rightarrow (0,-6)$; corte con $X$ en $2x - 6 = 0 \Rightarrow x = 3 \Rightarrow (3,0)$.
          \\[2pt]
          b) $F(x) = -\dfrac{3}{2}x + 6$: corte con $Y$ en $F(0) = 6 \Rightarrow (0,6)$; corte con $X$ en $-\dfrac{3}{2}x + 6 = 0 \Rightarrow x = 4 \Rightarrow (4,0)$.

    \item Corte con $Y$: $F(0) = 8 \Rightarrow (0,8)$. Corte con $X$: $\dfrac{4}{5}x + 8 = 0 \Rightarrow x = -10 \Rightarrow (-10, 0)$.
          Como $b = 8 \neq 0$, es una función \textbf{afín} (creciente, $m = \frac{4}{5} > 0$).
  \end{enumerate}

  \newpage

  \subsection*{Sección 5: Análisis de la Pendiente}

  \begin{enumerate}[leftmargin=*, label=\textbf{\arabic*.}]
    \item a) $m = \dfrac{11 - 1}{3 - (-2)} = \dfrac{10}{5} = 2$ \quad
          b) $m = \dfrac{-2 - 7}{4 - 1} = \dfrac{-9}{3} = -3$ \quad
          c) $m = \dfrac{-4 - 5}{3 - 3} = \dfrac{-9}{0}$: \textbf{no definida} (recta vertical; no es función).

    \item a) decreciente ($m = -3$) \quad b) constante ($m = 0$) \quad c) decreciente ($m = -1$) \quad d) creciente ($m = \frac{5}{2}$)

    \item $F(x) = -4x + b$ con $F(2) = -1$: $-8 + b = -1 \Rightarrow b = 7$. \quad $F(x) = -4x + 7$.

    \item $m = \dfrac{-3 - 6}{(k+3) - k} = \dfrac{-9}{3} = -3$ (la $k$ se cancela).
  \end{enumerate}

  \newpage

  \subsection*{Sección 6: Problemas de Aplicación}

  \begin{enumerate}[leftmargin=*, label=\textbf{\arabic*.}]
    \item Precio por kilo: $m = \dfrac{5\,600}{4} = 1\,400$. Función: $F(x) = 1\,400x$.
          $F(7) = 1\,400 \cdot 7 = \$9\,800$.

    \item $F(x) = 700x + 2\,000$ (afín). \quad $F(9) = 700(9) + 2\,000 = 6\,300 + 2\,000 = \$8\,300$.

    \item a) $F(x) = 4\,000x + 12\,000$.
          \begin{enumerate}[label=\alph*)]
            \setcounter{enumii}{1}
            \item $F(8) = 4\,000(8) + 12\,000 = \$44\,000$.
            \item $4\,000x + 12\,000 = 60\,000 \Rightarrow 4\,000x = 48\,000 \Rightarrow x = 12$ meses.
          \end{enumerate}

    \item $F(v) = 0{,}08v + 300\,000$. \quad $0{,}08v + 300\,000 = 460\,000 \Rightarrow 0{,}08v = 160\,000 \Rightarrow v = \$2\,000\,000$.
  \end{enumerate}

\fi
\end{document}
